\chapter{Noise Robustness Across Molecular Representations and Model Architectures}\label{chap:robustness}
% TODO: Chapter 3 update pending. When it lands, search for `TODO: Ch3 update' in Chapters 1, 4 and 7 and re-check each quoted result. Also re-check the Chapter 3 Additional files and the abstract.
% [notes moved to thesis_notes_moved.md]

% TODO: review, needs manual review
% [CODE INPUT: commented out -- the Chapter 3 paper was not submitted, so it is not referenced]
% \textit{Preamble.} The work in this chapter is adapted from Punt, A., Hanser, T., Werner, S., and Morris, G.M., in preparation for the \textit{Journal of Cheminformatics}. Most of the figures, tables, and text in this chapter are therefore either identical or similar to the content of that manuscript. A signed Statement of Authorship for this jointly authored work is included in the front matter.

\section{Introduction}
% [notes moved to thesis_notes_moved.md]


While there have been many investigations into the effect of noise on molecular property prediction, direction comparisons between how different molecular representations responds to label noise remains limited. The choice of molecular representation certainly affects predictive performance. \citet{Deng2023} evaluated 62,820 models and found that ECFPs frequently outperform learned graph embeddings across models. Whether these conclusions translate to noise robustness is our question. Between different types of label noise, \citet{Heid2023} found no difference in model performance between error distributions with the same mean and standard deviation. We expand on this by considering alternative patterns of noise, namely noise concentrated in structural groups or on a sparse subset of samples.

% TODO: review, needs manual review
This chapter addresses the first four questions set out in Chapter~\ref{chap:intro}, namely how representation and model architecture divide the variance in predictive performance and in noise robustness, what probabilistic models add, whether the type of label noise matters, and whether uncertainty estimates track label noise.

\section{Methods}\label{sec:robustness:methods}
% ---- NOTES: Methods ----
% Takeaway: the study crosses four datasets, six representations, 20 model configurations, seven noise conditions and seven noise levels.
% Paragraph 1 (only if the section gets a lead-in paragraph):
% - Four datasets: QM9 (HOMO--LUMO gap) and three assay datasets (logD, Caco-2, hERG K$_i$).
% - Six representations, 20 model configurations, seven noise conditions, seven noise levels.
% - The 20 configurations are thirteen base models and seven variant models. Name both terms here; the Models subsection defines them.
% - The noise conditions are defined in Section~\ref{sec:noise_injection}. Hyperparameters are in Additional file~\ref{app:noise:sec:hyperparameters_default}.
% ---- END NOTES ----

% CODE INPUT
We used four datasets, namely QM9 with the HOMO--LUMO gap as the prediction target and three assay datasets for LogD, Caco-2 and hERG $K_i$. On each dataset we combined six representations with 20 model configurations, which comprise thirteen base models and seven variant models. Every pairing of representation and model was exposed to seven noise conditions (Section~\ref{sec:noise_injection}) at seven noise levels, and the hyperparameters of every model are given in Additional file~\ref{app:noise:sec:hyperparameters_default}.



\subsection{Datasets}

Although we are investigating the effects of label noise stemming from experimental settings, for the purposes of testing we inject controlled amounts of artificial noise into clean labels. Data were extracted from the QM9 molecular property dataset distributed by PyTorch Geometric \citep{Ramakrishnan2014}, which contains 130,831 small organic molecules with pre-computed quantum-mechanical properties obtained from density functional theory calculations at the B3LYP/6-31G(2df,p) level of theory. We selected the HOMO--LUMO gap as the prediction target in QM9, as it captures a molecule's electronic excitability, charge transfer capability, and chemical stability \citep{Fediai2023}. We assumed that QM9 labels were free of noise, notwithstanding approximations in the level of theory. Each experiment was replicated ten times with a distinct random seed and a randomly selected subset of $N = 10{,}000$ molecules, split 80/10/10 into 8,000 training, 1,000 validation and 1,000 test molecules. All models were fitted on those 8,000 training molecules, except the Gaussian processes, which were fitted on a random subset of 5,000 of them; the validation and test molecules were unchanged.

We repeated these experiments on three molecular property datasets with experimentally determined endpoints. From the OpenADMET initiative \citep{openadmet}, we used LogD and Caco-2 efflux permeability, with $N = 5{,}039$ and $N = 2{,}161$ respectively. We also selected hERG $K_i$ data from ChEMBL \citep{Zdrazil2023}, following a protocol inspired by \citet{landrum2024}: filter for binding assays, deduplicate by median pChEMBL value, and remove compounds with an inter-assay standard deviation greater than 1.0 log unit, resulting in $N = 1{,}415$ compounds. Each of the three was partitioned into five folds by grouped $k$-fold cross-validation, and each fold's training block gave up a further grouped fifth as a validation set for early stopping. A model therefore fits between 3,194 and 3,364 of the 5,039 LogD molecules, depending on the fold. Unlike QM9's ten replicates, which are ten independent draws from one pool, each assay dataset is a single fixed set in which every molecule is tested exactly once and nothing is redrawn.

All splits and cross-validation folds were grouped by the chirality-blind Bemis--Murcko scaffold, implemented such that each distinct acyclic molecule forms its own group, and groups are filled in random order. Prior to modeling, all molecular structures were sanitized and canonicalized using \texttt{RDKit} to remove invalid valence states and standardize atom and bond typing \citep{rdkit}. Prediction targets were standardized to zero mean and unit variance. Experiments that track uncertainty were run once rather than in replicates.

\subsection{Representations}

In this study, we evaluated six molecular representations: ECFP4, Avalon, Sort \& Slice, PDV, ChemBERTa and Molecular Hypergraph Grammar GNNs (MHG-GNNs). The same six were used on QM9 and on the three assay datasets. ECFP4 is an ECFP with radius $r = 2$ with $d = 2048$ bits \citep{rogers2010, rdkit}. Sort \& Slice takes the same substructures as ECFP4, but it ranks them by prevalence and selects the top $L$ ($L=1024$) \citep{sns}. We fitted the Sort \& Slice vocabulary on the training molecules alone, refitting it for each QM9 replicate and for each cross-validation fold.  Avalon is a 2,048-bit fingerprint, generated with the Avalon toolkit through RDKit \citep{avalon, rdkit}. The PDV representation is a vector of 200 physicochemical descriptors computed with RDKit's \texttt{MolecularDescriptorCalculator} \citep{rdkit, Cherkasov2014}; the 200 are listed in Additional file~\ref{app:noise:sec:pdv_descriptors}, and the list is fixed in our code rather than taken from RDKit's default descriptor set, which grows between releases. ChemBERTa embeddings are the mean over the non-padding token embeddings of the \texttt{DeepChem/ChemBERTa-77M-MTR} checkpoint, producing 384 features \citep{Ahmad2022}. MHG-GNNs are a GIN-based autoencoder pre-trained on 1.34 million PubChem molecules using a $\beta$-variational autoencoder loss, producing 1024-dimensional embeddings through iterative message passing and found in the published \texttt{mhggnn\_pretrained\_model\_0724\_2023} checkpoint \citep{kishimoto2023}. We z-score normalized PDV, ChemBERTa and MHG-GNN per feature, using the training molecules' mean and standard deviation; ECFP4, Avalon and Sort \& Slice reach the model as built.

\subsection{Models}
% [notes moved to thesis_notes_moved.md]

% TODO: review, needs manual review
To understand model behavior in the presence of label noise, we fitted models from several families, namely tree ensembles (random forests, quantile regression forests, XGBoost, LightGBM and NGBoost), a support vector machine, Gaussian processes and neural networks. Several of these models are variants of one another that differ only in the probabilistic components they add, and the hyperparameters of all models described below are listed in Additional file~\ref{app:noise:sec:hyperparameters_default}.

While we used 100 trees for RF, we observed a noisy quantile estimate with 100 trees for QRF, prompting us to use 300 instead. To observe the effect of the quantile mechanism, we then fit a second RF at 300 trees so the two models were idential in all other settings.

NGBoost relies on scoring off validation data, which differs by dataset. On QM9 the validation split is already held out of every fit. However, on the three assay datasets the validation set is a scaffold-grouped fifth of the rows NGBoost is handed, which already exclude the validation molecules.

We used an RBF with $C = 1.0$ and $\gamma = \texttt{scale}$ as the kernel, as it accomodates continuous vectors.

We fitted three variants of GPs, each with a constant mean and a scaled kernel. The first uses an RBF kernel and runs on all six representations. The second uses the Tanimoto kernel \citep{Ralaivola2005, moss2020} and was run on ECFP4 alone. The third is a heteroscedastic variant in which the RBF process is given a second network that predicts the observation noise of each molecule. As such, we used a standard GP and only fit on at most $5{,}000$ training molecules, drawn by a seeded sample. Due to the size of the datasets, this cap only affected QM9. We implemented these GP models in \texttt{GPyTorch} \citep{gpytorch}, with the Tanimoto kernel taken from the \texttt{Gauche} framework \citep{gauche}.

While preliminary tests indicated that the Tanimoto kernel provides GP with substantial benefits over other kernels, for SVM, the RBF kernel was sufficient for all types of data.

% TODO: review, needs manual review
We tested two deterministic, fully connected feed-forward architectures, both implemented in PyTorch \citep{pytorchGeometric}. NN-$\alpha$ has two hidden layers of 128 and 64 units with dropout after each hidden layer, and NN-$\beta$ has two hidden layers of 128 units with dropout applied before the output layer. Both were trained with the Adam optimizer and early stopping on the validation loss. Bayesian transformations converted these networks into Bayesian neural networks (BNNs) with ReLU activations and dropout regularization ($p=0.2$), and Additional file~\ref{app:noise:sec:hyperparameters_default} gives the full widths, dropout fractions, learning rates and batch size of every network.

To get these networks producing uncertainty estimates, we applied four probabilistic transformations to each architecture, producing eight probabilistic networks. The first replaces every linear layer with a Bayesian layer carrying a Gaussian prior $\mathcal{N}(0, 0.1^2)$ on its weights, producing BNN-$\alpha$ and BNN-$\beta$. They are trained on an evidence lower bound whose  Kullback–Leibler (KL) divergence term is divided by the number of training molecules. While this is a straightforward approach, it is computationally expensive. We explored alternative techniques like our second variant, where we add an output head to the standard NN and fit the negative log likelihood of a Gaussian with that variance \citep{kendall2017}. The third, Variational Bayesian Last Layer (VBLL), replaces every linear layer with a layer that produces mean-field variational posteriors $q(\mathbf{W}) = \mathcal{N}(\boldsymbol{\mu}_W, \text{diag}(\boldsymbol{\sigma}_W^2))$ over the last-layer weights and trains by maximizing the evidence lower bound. This can be thought of as minimizing the reconstruction loss plus KL divergence $D_{\text{KL}}(q(\mathbf{W}) \| p(\mathbf{W}))$ from a standard normal prior, scaled by $1/N$ \citep{Harrison2024}. Because only its output layer keeps a learned observation-noise parameter, VBLL networks report only epistemic uncertainty. The fourth makes observation noise a function of the input. All eight probabilistic networks estimate their predictive distributions from 100 stochastic forward passes. Each network and its Bayesian counterpart were tuned to the same set of parameters so that NN-$\alpha$ against BNN-$\alpha$ and NN-$\beta$ against BNN-$\beta$ differ by the Bayesian layers alone.

\subsection{Noise Conditions and Levels}
% [notes moved to thesis_notes_moved.md]
% TODO: review, needs manual review
The seven noise conditions used in this chapter are defined in Section~\ref{sec:noise_injection}.

We ran experiments across seven noise levels, $0$, $0.2$, $0.3$, $0.5$, $0.75$, $1.0$ and $1.5$. We added the noise to the raw label and standardized the target afterward, using the clean training mean and spread. For censoring, the level is the fraction of labels clipped, at $0$, $0.1$, $0.2$, $0.25$, $0.3$, $0.4$ and $0.5$ \citep{Svensson2025}.

We ran Gaussian, grouped-wider and grouped-shifted noise on every model and representation pairing. The
remaining four conditions ran on selected pairs. 


\subsection{QM9 Data Pipeline}
% [notes moved to thesis_notes_moved.md]

% TODO: review, needs manual review
QM9 \citep{Ramakrishnan2014}, loaded through PyTorch Geometric, was processed by a separate pipeline
in which Python code loads and splits the data, computes every representation, and trains and scores
the models. Between featurization and training, a program written in Rust injects the noise,
standardizes the labels and writes the files from which the models are trained.

% TODO: review, needs manual review
Each replicate used a new scaffold split, drawn as described in the Datasets subsection above. Each
acyclic molecule forms its own group, identified by its canonical SMILES, and no group is ever divided
between two splits. The same grouping defines the scaffold families targeted by the grouped noise
conditions.

% TODO: review, needs manual review
The six representations described above are computed in Python. Binary fingerprints are stored in
bit-packed form, Sort \& Slice counts as 16-bit integers, and the descriptors and learned embeddings
as 32-bit floating-point numbers. The run halts if any molecule has an ECFP4 or Avalon fingerprint
that is entirely zero, so that a molecule that failed featurization cannot enter training. The
pipeline can also build one-hot encoded SMILES, although these were not used in this thesis.

% TODO: review, needs manual review
Conventional approaches load entire datasets into memory from comma-separated values (CSV) or Parquet files, which creates memory bottlenecks at scale and repeats the cost of deserialization in every experimental replicate. Our pipeline instead serializes the preprocessed data into custom memory-mapped binary files that are read on demand, so that memory use stays low and constant regardless of dataset size (Figure~\ref{fig:fig_mmap}). Python writes one file per split, the Rust program reads each file and writes a new one in which the processed labels replace the raw ones, and Python then reads the new files back for training.

% TODO: review, needs manual review
Each memory-mapped file stores a sequence of fixed-format entries. An entry holds the isomeric and canonical SMILES strings in UTF-8 with a length prefix, the raw target value as a 4-byte float, the requested representations as fixed-width blocks, and the processed target value, which includes the noise and normalization, as a 4-byte float. The blocks occupy 2,048 bytes for Sort \& Slice (1,024 16-bit integers), 800 bytes for PDV (200 32-bit floats), 1,536 bytes for ChemBERTa (384 32-bit floats), 4,096 bytes for MHG-GNN (1,024 32-bit floats), and 256 bytes each for Avalon and ECFP4. Only the requested representations are written, so the byte offset of every field depends on the list of representations. This format permits zero-copy memory mapping, in which Python and Rust access the same disk-backed memory pages without deserialization and memory use stays constant regardless of dataset size.

% TODO: review, needs manual review
The Rust program draws the noise for each split separately, following the order in which that
split's records are stored. The noise parameters are set on the training split, since the amount of
noise, the censoring limit and the affected scaffold families are all computed from the clean
training labels. Validation labels receive an independent draw, from a seed derived from the
training seed, with the same amount of noise in label units, the same limit and the same affected
families, and the program verifies that validation received as much noise as training. Test labels
are never corrupted, although every split records the per-molecule noise scale its region would
have received, so that held-out molecules can be scored against the pattern applied in training.

% TODO: review, needs manual review
For every molecule in every split, the Rust program records its canonical SMILES, clean label,
injected noise, noise scale, the noise pattern applied to its region, noisy label and the
standardized value written to file. It also writes a summary of each run that lists the parameters
of the noise condition and the amount of noise delivered. The run halts if the record at any row
index holds a different molecule when it is written than when its noise was drawn. A separate
self-test mode applies the same checks on the amount and targeting of noise to a label file,
without running the rest of the pipeline.

% TODO: review, needs manual review
\paragraph{Model training.} The pipeline trains the models described in the Models subsection above,
with hyperparameters taken from package defaults, tuned with Optuna \citep{akiba2019optuna}, or read
from a JSON configuration file. For models that return uncertainty, the pipeline can also score the
training molecules out of fold, so that the uncertainty of each training molecule comes from a model
that never saw its corrupted label.

\begin{figure}[htbp]
    \centering
    \includegraphics[width=0.95\textwidth]{figure_mmap_structure.png}
    \caption[In-memory and memory-mapped data loading.]{Schematic comparison of in-memory and memory-mapped data loading. The top left shows conventional loading, in which the entire CSV or Parquet dataset is read into random-access memory (RAM). The bottom left shows memory-mapped loading, in which each molecule is written as one line of a binary file saved on disk and lines are read into RAM on demand. The table on the right lists the fields of one line, with whether each field is required, its data type and its size in bytes.}
    \label{fig:fig_mmap}
\end{figure}


\subsection{Uncertainty}
% [notes moved to thesis_notes_moved.md]

Up to this point we have been concerned with how noise affects the predictive accuracy of QSAR models. However, label noise should also affect the uncertainty these models report. Some models return only a point prediction, while others also return an estimate of uncertainty, either as a predicted distribution or as the spread across an ensemble. Furthermore, some report a single value of uncertainty while others report one that can be decomposed into an aleatoric and an epistemic component \citep{kendall2017}. The aleatoric component describes underlying noise the model associates with the data, while the epistemic component represents the uncertainty within the model's fit. We sought to determine if, as expected, the aleatoric component rises with injected label noise while the epistemic component does not.

For the BNN and VBLL transformations, we took the two components over $T = 100$ stochastic forward
passes \citep{kendall2017}:
\begin{equation}
\hat{u}^2_{\text{ale}}(x) = \frac{1}{T}\sum_{t=1}^{T}\hat{v}_t(x),
\qquad
\hat{u}^2_{\text{epi}}(x) = \frac{1}{T}\sum_{t=1}^{T}\big(\hat{\mu}_t(x) - \bar{\mu}(x)\big)^2 ,
\label{eq:sampling_split}
\end{equation}
where $\hat{\mu}_t(x)$ and $\hat{v}_t(x)$ are the mean and the variance returned for molecule $x$ on
pass $t$, and $\bar{\mu}(x)$ is the mean over passes. Here $\hat{v}_t(x)$ is the predicted observation
noise variance. It varies per molecule in the variance-head and heteroscedastic variants, and is one
number for the whole fit in the VBLL transformation. A full-BNN without a variance head returns no
variance, so its uncertainty is epistemic alone.

The other models follow the same split. For GPs, the epistemic component is the posterior variance of the latent function and the aleatoric component is the learned observation noise \citep{Rasmussen2005}. All variants of the GP produce one number for the whole fit except heteroscedastic. For QRF models, each tree holds multiple training molecules. So for a given molecule, the aleatoric component is the variance the labels the molecule's given leaf. The total aleatoric uncertainty are those variances averaged over the trees. The epistemic component is the variance of the leaf means across the trees. NGBoost predicts one variance for each molecule from a single fit, which is aleatoric alone \citep{Duan2020}. In every model, the total uncertainty is the sum of the the aleatoric and epistemic components.

No held-out label carries injected noise, so every comparison against the injected noise was made on
training molecules. Each training set was divided into five folds grouped on Murcko scaffolds, and each
molecule was scored by a refit that excluded its own fold. For NGBoost on QM9, three of the five folds
were scored to limit its cost. No uncertainty was calibrated after fitting.


\subsection{Evaluation}
% [notes moved to thesis_notes_moved.md]
We compute statistics on the held-out molecules, for QM9 the test split and for three assay datasets the held-out fold of the scaffold cross-validation. Different combinations of dataset, model, representation, noise condition, and noise level are examined in each scenario. We used the coefficient of determination ($R^2$) throughout. For molecule $i$, let $y_i$ be the clean label, $\hat{y}_i$ the prediction, $\epsilon_i$ the noise injected into its label and $u_i$ the predicted uncertainty.

To evaluate the effect of label noise, we tracked predictive performance retention under increasing artificial noise. For each QSAR pairing we recorded $R^2(\tau)$ at each noise level $\tau$, and normalized it to obtain the retention ratio $R^2(\tau)/R^2(0)$. To quantify robustness, we used the normalized area under this retention curve, \aucnorm (Equation~\ref{eq:auc_norm}), evaluated using the trapezoidal rule over the seven noise levels. While we tested other metrics of robustness, we found \aucnorm to be less directly correlated to clean-label $R^2$ predictions. For censoring, we sweep over a set of fractions to clip the labels prior to normalization. 

An \aucnorm value near 1 indicates that the QSAR model in question has retained predictive performance with artificial noise, meaning it is more robust to noise. Since we normalize the curve by its own clean-label predictive performance, a model with a low clean $R^2$ has less to lose, so \aucnorm is reported beside clean $R^2$. We excluded replicates with a clean $R^2 < 0.3$ from the robustness analysis (Additional file~\ref{app:noise:sec:excluded_configurations}), as performance retention ratios become unstable when the clean-label denominator approaches zero. On QM9 one \aucnorm is computed per replicate, and we report the median over the ten replicates. On the three assay datasets one value is computed per scaffold fold, and the median over the five folds is reported. 

To compare the relative effect of the thirteen base models across datasets and noise conditions, we compared their rankings using Spearman's rank correlation coefficient ($\rho$) between clean $R^2$ and \aucnorm. We compared probabilistic models  with its deterministic counterpart using a two-sided Wilcoxon signed-rank test paired on the replicate, at $\alpha = 0.05$, separately for each representation and noise condition. This test was run only on QM9, as with five scaffold folds the
smallest attainable $p$ is 0.0625.

We conducted a separate two-way analysis of variance (ANOVA) for each noise condition to identify the
relative contributions of molecular representation and model architecture to both predictive
performance and noise robustness. For a given metric $y$, either $R^2$ at a noise level of 1.0 or
\aucnorm:
\begin{equation}
y_{ijr} = \mu + \alpha_i + \beta_j + (\alpha\beta)_{ij} + e_{ijr},
\end{equation}
where $\alpha_i$ represents the effect of model architecture $i$, $\beta_j$ represents the effect of
molecular representation $j$, $(\alpha\beta)_{ij}$ is the interaction term, and $e_{ijr}$ is the
residual for replicate $r$. The proportion of variance explained by each term was calculated as the
effect size
\begin{equation}
\eta^2_{\text{term}} = \frac{SS_{\text{term}}}{SS_{\text{total}}},
\end{equation}
using Type I (sequential) sums of squares, entered in the order model architecture, molecular
representation, interaction. The ten QM9 replicates provide the residual. The uncertainty on each share is the range of $\eta^2$ when the decomposition is repeated with each replicate left out in turn. 

ANOVA relies on the assumption that observations both within and across groups are independence. Across the whole set of models and representations in this study, some pairs produced similar predictions that are not independent. This causes the residual degrees of freedom to increase, while the residual sum of squares barely moves. To handle this issue, we computed pairwise Spearman rank correlations between all models and between all representations across noise strategies. Models whose profiles correlated at $\rho > 0.99$ with another model were excluded, including QRFs, which were very similar to RF, as well as different BNN variants. 

Across representations, Sort \& Slice correlated at $\rho > 0.90$ with the ECFP4 fingerprint and was excluded to avoid inflating representation degrees of freedom, as both encode overlapping circular substructure information. We also computed intraclass correlation coefficients ICC(1,1) for all model pairs. ICC(1,1) measures the proportion of total variance attributable to between-subject differences, with values near 1.0 indicating that two models rank configurations almost identically. The full redundancy and ICC tables are reported in Additional files~\ref{app:noise:sec:model_redundancy}--\ref{app:noise:sec:model_icc}.

To assess the interaction between uncertainty estimation and label noise, we posed three questions. First, whether uncertainty rise with label noise, as measured by the mean $u_i$ over held-out molecules at the highest noise level
divided by its value with no added noise, in each replicate or fold. The second pertains to models who report separate aleatoric and epistemic components. We divided the increase of mean epistemic uncertainty between the highest level of noise and clean labels by the mean of the aleatoric component to determine how each responds to label noise. A ratio near 0 means the added noise was attributed to the labels while a ratio of 1 or above means the epistemic component grew as much or more. We also compared the mean aleatoric component with the standard deviation of the injected noise. The third question we had was whether per-sample uncertainty values track injected label noise, as measured by the Spearman correlation between $u_i$ and $|\epsilon_i|$ over out-of-fold training molecules, per fold and per noise level. We also report the Spearman correlation between $u_i$ and the absolute error of the prediction. Finally, we asked whether dividing a molecule's out-of-fold absolute error by $u_i$ ranks the corrupted labels better than the error alone, and report the difference between the two Spearman correlations with $|\epsilon_i|$.


\section{Variance Decomposition by Representation and Architecture}
% [notes moved to thesis_notes_moved.md]


Considering the two main components of a QSAR model, the molecular representation and the model architecture, we sought to understand how each contributes to predictive performance and noise robustness. ANOVA decomposition across six noise strategies on models to predict HOMO--LUMO gap from QM9 demonstrates that for predictive performance (R$^2$ at $\sigma = 0.3$), the model--representation interaction term is the  largest source of variance across all six strategies, followed by model architecture and representation (Table~\ref{tab:anova_decomposition}; Figure~\ref{fig:anova_decomposition}a). However, for noise robustness as measured by the noise degradation slope (NDS), this ordering shifts. The choice of model architecture is instead the largest source of variance, while the interaction term drops to second and representation third. This observed inversion makes sense because the molecular representation dictates the chemical information the model receives, while the architecture determines how efficiently that information is used. The interaction term reflects the idea that these two are not independent. In the context of label noise robustness, representation has less of an impact since the features are not corrupted by noise. Meanwhile, the model's regularization, ensembling, and uncertainty handling have a stronger impact on how it handles noise. 

The impact varied by type of noise, represented by noise strategy (Figure~\ref{fig:anova_decomposition}b). The strategies which degraded in performance the most, threshold and value-proportional, both corrupted label regions, forcing models to rely on their inductive biases. This translates to the model architecture explaining a significant portion of the variance. On the other hand, outlier noise only affects statistical outliers. For outlier noise, model and interaction effects were both smaller.

\begin{table}[htbp]
\centering
\caption[ANOVA variance decomposition by noise strategy on QM9.]{Two-way analysis of variance (ANOVA) decomposition by noise strategy on the QM9 HOMO--LUMO gap. Each entry is the share of variance $\eta^2$ (\%) explained by the model architecture (Model), the molecular representation (Rep), the model--representation interaction (Inter.) or the residual (Resid.). The performance columns use $R^2$ at $\sigma = 0.3$ as the outcome, and the robustness columns use the noise degradation slope (NDS). Value-prop.\ denotes value-proportional noise. Bold marks the largest share within the performance columns and within the robustness columns of each row.}
\label{tab:anova_decomposition}
\small
\begin{adjustbox}{max width=\textwidth}
\begin{tabular}{lcccccccc}
\toprule
& \multicolumn{4}{c}{\textbf{Performance ($\eta^2$, \%)}} & \multicolumn{4}{c}{\textbf{Robustness ($\eta^2$, \%)}} \\
\cmidrule(lr){2-5} \cmidrule(lr){6-9}
\textbf{Strategy} & \textbf{Model} & \textbf{Rep} & \textbf{Inter.} & \textbf{Resid.} & \textbf{Model} & \textbf{Rep} & \textbf{Inter.} & \textbf{Resid.} \\
\midrule
Gaussian        & 24.6 & 22.8 & \textbf{48.1} &  4.5 & \textbf{48.7} &  8.2 & 11.9 & 31.2 \\
Quantile        & 25.1 & 22.4 & \textbf{46.9} &  5.6 & \textbf{47.9} &  6.8 & 14.9 & 30.4 \\
Threshold       & 25.0 & 21.2 & \textbf{48.0} &  5.8 & \textbf{48.0} & 10.1 & 28.0 & 13.9 \\
Heteroscedastic & 27.9 & 23.5 & \textbf{46.2} &  2.5 & \textbf{37.0} &  4.5 & 17.5 & 41.0 \\
Value-prop.     & 27.8 & 22.0 & \textbf{46.5} &  3.8 & \textbf{52.8} &  6.9 & 22.2 & 18.2 \\
Outlier         & 25.3 & 22.0 & \textbf{49.2} &  3.5 & 12.7 &  3.2 &  4.7 & \textbf{79.3} \\
\bottomrule
\end{tabular}
\end{adjustbox}
\end{table}

\begin{figure}[htbp]
    \centering
    \includegraphics[width=\textwidth]{fig2_anova_decomposition.png}
    \caption[ANOVA variance decomposition of performance and robustness on QM9.]{Analysis of variance (ANOVA) decomposition of performance and robustness by noise strategy on the QM9 HOMO--LUMO gap. Each bar is the share of variance $\eta^2$ (\%) explained by model architecture, molecular representation or their interaction (indicated by colors). The top panel (a) decomposes predictive performance, measured as $R^2$ at $\sigma = 0.3$, and the bottom panel (b) decomposes noise robustness, measured as the noise degradation slope (NDS). The residual share is not plotted and is given in Table~\ref{tab:anova_decomposition}.}
    \label{fig:anova_decomposition}
\end{figure}

The interaction between model architecture and representation is seen in Figure~\ref{fig:interaction}. While some model architectures, like SVM and full BNNs maintain consistent noise robustness scores across all representations, others like NN-$\beta$ display greater variation. Model robustness on one representation strongly predicts robustness on another (Spearman $\rho = 0.73$, $p = 0.007$ for ECFP4 vs PDV; 11 models). However, the scatter around the diagonal reflects the interaction term identified by the ANOVA. 

\begin{figure}[htbp]
\centering
\includegraphics[width=\textwidth]{fig_interaction.png}
\caption[Model--representation interaction in noise robustness on QM9.]{Model--representation interaction in noise robustness on the QM9 HOMO--LUMO gap under Gaussian noise. The top panel (a) shows the noise degradation slope (NDS) of each model on each representation, where more negative values (darker orange) indicate faster loss of performance. Black cells marked N/A are configurations excluded because their clean $R^2$ was below 0.6 or because the kernel was incompatible with the representation, as for the Gaussian process (GP) with a Tanimoto kernel, which is defined on binary vectors and cannot be applied to the physicochemical descriptor vector (PDV). The bottom panel (b) plots the NDS on PDV against the NDS on the extended-connectivity fingerprint (ECFP4) for each model, with the dashed line marking equal values and the Spearman rank correlation $\rho$ given in the panel.}
\label{fig:interaction}
\end{figure}

We filtered out model--representation configurations with baseline predictive performance R$^2 < 0.6$, excluding 66 configurations where poor clean-data performance would produce misleadingly shallow degradation slopes (Additional file~\ref{app:noise:sec:excluded_configurations}). We evaluated all model architectures under increasing label noise on the PDV representation across six strategies, considering both predictive performance as measured by $R^2$ as well as the NDS. We observe that NDS varies substantially across architectures and noise strategies, as seen in Figure~\ref{fig:global_overview}. 

\begin{figure}[htbp]
\centering
\includegraphics[width=\textwidth]{fig1_global_overview.png}
\caption[Noise robustness of all models on PDV for QM9.]{Overview of noise robustness on the QM9 HOMO--LUMO gap with the physicochemical descriptor vector (PDV) representation. The top panel (a) shows $R^2$ against the noise level $\sigma$ under Gaussian noise for seven representative models. The bottom panel (b) shows the noise degradation slope (NDS) of every model under each of six noise strategies, where more negative values (orange) indicate faster loss of performance and values closer to zero (purple) indicate slower loss. The Gaussian process (GP) with a Tanimoto kernel is absent because its kernel is incompatible with PDV (see Section~\ref{sec2}), and the GP with a radial basis function kernel, GP (RBF), is shown instead.}
\label{fig:global_overview}
\end{figure}

\section{Probabilistic and Deterministic Model Variants}
% [notes moved to thesis_notes_moved.md]

To further explore the impact of model architecture on noise robustness, we examined variants of the same core architecture. Specifically, Bayesian transformations were used to explore the impact of the Bayesian framework on noise. By comparing models with the same core architecture, we observe that with full Bayesian transformations, representation explains less than $11\%$ of NDS variance ($p > 0.25$) while VBLL on the same NNs is representation-dependent. On MHG-GNN and Mol2vec representations, VBLL models were unstable, such that all 24 VBLL $\times$ \{MHG-GNN, Mol2vec\} configurations were excluded due to baseline R$^2 < 0.6$. Further work could involve more thoroughly-tuning these models and exploring different core architectures to adapt to these embedding-style inputs. 

In general, the full Bayesian transformations reduced the degradation in performance with Gaussian noise the most, as seen in Figure~\ref{fig:nn_family_comparison}. Both full BNN and VBLL transformations significantly improved robustness for both NN-$\alpha$ and NN-$\beta$ (Table~\ref{tab:wilcoxon_bnn}). Full transformations are more computationally expensive, making VBLL a potentially stronger approach to handling noise with NNs when working with larger data sets or modeling architectures. The choice of transformation varies by representation, as seen in Additional file~\ref{app:noise:sec:bayesian_transformation_tests}, with SMILES receiving the largest benefits from Bayesian transformations. QRF was significantly less robust than RF (Table~\ref{tab:wilcoxon_bnn}), suggesting that quantile regression adds complexity and flexibility which may result in overfitting.

\begin{table}[htbp]
\centering
\caption[Wilcoxon signed-rank tests of probabilistic variants on QM9.]{Wilcoxon signed-rank tests comparing each deterministic base model with its Bayesian or probabilistic variant on the QM9 HOMO--LUMO gap. $\Delta$NDS is the mean change in noise degradation slope (NDS) from the base model to the variant, where positive values mean that the variant is more robust (slope closer to zero). An asterisk in the Sig.\ column marks pairs with $p < 0.05$, and bold $\Delta$NDS values mark the largest positive change and the only negative change. NN = neural network, BNN = Bayesian neural network, VBLL = Variational Bayesian Last Layer, RF = random forest, QRF = quantile regression forest.}
\label{tab:wilcoxon_bnn}
\small
\begin{adjustbox}{max width=\textwidth}
\begin{tabular}{llrrl}
\toprule
\textbf{Base} & \textbf{Variant} & \textbf{$\Delta$ NDS} & \textbf{$p$-value} & \textbf{Sig.} \\
\midrule
NN-$\alpha$ & BNN-$\alpha$ & $+0.056$ & $2.9 \times 10^{-10}$ & \textbf{*} \\
NN-$\alpha$ & VBLL-$\alpha$ & $+0.061$ & $1.2 \times 10^{-6}$ & \textbf{*} \\
NN-$\beta$ & BNN-$\beta$ & $+0.096$ & $2.9 \times 10^{-11}$ & \textbf{*} \\
NN-$\beta$ & VBLL-$\beta$ & $\mathbf{+0.124}$ & $1.2 \times 10^{-7}$ & \textbf{*} \\
RF  & QRF  & $\mathbf{-0.022}$ & $1.6 \times 10^{-10}$ & \textbf{*} \\
\bottomrule
\end{tabular}
\end{adjustbox}
\end{table}

\begin{figure}[htbp]
    \centering
    \includegraphics[width=\textwidth]{fig_nn_family_comparison.png}
    \caption[Neural network families and random forests under Gaussian noise on QM9.]{Predictive performance of neural network (NN) families and random forests under Gaussian noise on the QM9 HOMO--LUMO gap ($N = 10{,}000$) with the physicochemical descriptor vector (PDV) representation. Each panel shows $R^2$ against the noise level $\sigma$. The top panel (a) shows NN-$\alpha$ with its full Bayesian (BNN-$\alpha$) and Variational Bayesian Last Layer (VBLL-$\alpha$) variants, the middle panel (b) shows the corresponding NN-$\beta$ family, and the bottom panel (c) compares the random forest (RF) with the quantile regression forest (QRF).}
    \label{fig:nn_family_comparison}
\end{figure}

\section{Model Robustness Rankings across Noise Strategies}
% [notes moved to thesis_notes_moved.md]

NGBoost and SVM showed the smallest degradation in predictive performance across all strategies, while NN-$\beta$ showed the steepest (Table~\ref{tab:nds_ranking}). When sorting models by NDS, their ranks across noise strategies were consistent: Kendall's $W = 0.92$ ($p = 2.7 \times 10^{-8}$; 11 models, 6 strategies), indicating that a model's relative robustness is largely strategy-independent. The same ranking pattern holds on the ECFP4 fingerprint (Additional file~\ref{app:noise:sec:ecfp4_overview}). Robustness is also largely decoupled from baseline predictive performance: across models on PDV under Gaussian noise, NDS clusters near $-0.38$ regardless of baseline R$^2$ (Additional file~\ref{app:noise:sec:clean_r2_vs_robustness}), so a high-performing model is not automatically a more robust one.

\begin{table}[htbp]
\centering
\caption[Noise degradation slope by model on PDV for QM9.]{Noise degradation slope (NDS) of each model on the physicochemical descriptor vector (PDV) representation for the QM9 HOMO--LUMO gap, under each of six noise strategies and as the mean across them, with models ordered by the mean. A slope closer to zero indicates greater robustness. Gauss.\ = Gaussian, Quant.\ = quantile, Thresh.\ = threshold, Hetero.\ = heteroscedastic and Val.-P.\ = value-proportional noise. Bold model names mark the two models whose mean is closest to zero, and bold entries mark the value closest to zero in each column.}
\label{tab:nds_ranking}
\small
\begin{adjustbox}{max width=\textwidth}
\begin{tabular}{lccccccc}
\toprule
\textbf{Model} & \textbf{Gauss.} & \textbf{Outlier} & \textbf{Quant.} & \textbf{Thresh.} & \textbf{Hetero.} & \textbf{Val.-P.} & \textbf{Mean} \\
\midrule
\textbf{NGBoost}    & $\mathbf{-0.33}$ & $-0.12$ & $\mathbf{-0.29}$ & $\mathbf{-0.65}$ & $\mathbf{-0.18}$ & $\mathbf{-0.58}$ & $\mathbf{-0.36}$ \\
\textbf{SVM}        & $-0.36$ & $\mathbf{-0.11}$ & $-0.30$ & $-0.71$ & $-0.20$ & $-0.63$ & $-0.38$ \\
RF                  & $-0.36$ & $-0.12$ & $-0.30$ & $-0.71$ & $-0.20$ & $-0.63$ & $-0.39$ \\
LightGBM            & $-0.36$ & $-0.13$ & $-0.32$ & $-0.72$ & $-0.20$ & $-0.64$ & $-0.40$ \\
VBLL-$\alpha$       & $-0.38$ & $-0.16$ & $-0.35$ & $-0.68$ & $-0.22$ & $-0.60$ & $-0.40$ \\
XGBoost             & $-0.37$ & $-0.13$ & $-0.33$ & $-0.72$ & $-0.20$ & $-0.65$ & $-0.40$ \\
VBLL-$\beta$        & $-0.39$ & $-0.14$ & $-0.34$ & $-0.70$ & $-0.23$ & $-0.64$ & $-0.41$ \\
BNN-$\alpha$        & $-0.38$ & $-0.15$ & $-0.35$ & $-0.72$ & $-0.21$ & $-0.64$ & $-0.41$ \\
BNN-$\beta$         & $-0.38$ & $-0.15$ & $-0.36$ & $-0.72$ & $-0.21$ & $-0.65$ & $-0.41$ \\
NN-$\alpha$         & $-0.41$ & $-0.16$ & $-0.37$ & $-0.77$ & $-0.23$ & $-0.68$ & $-0.44$ \\
NN-$\beta$          & $-0.49$ & $-0.17$ & $-0.42$ & $-0.88$ & $-0.26$ & $-0.79$ & $-0.50$ \\
\bottomrule
\end{tabular}
\end{adjustbox}
\end{table}

The strategies differed not only in predictive performance severity, but also in their ability to discriminate between architectures. The spread of NDS across models differed by noise strategy, with outlier noise having a spread of 0.06 and threshold noise having 0.23, presumably exposing architectural differences that are masked by other types of noise. As seen in Table~\ref{tab:nds_ranking}, NGBoost was the most resistant to noise across strategies. With respect to the VBLL models, VBLL-$\alpha$ showed a similar pattern, outperforming all tree-based models except NGBoost under threshold and value-proportional noise despite ranking below them under Gaussian noise. This suggests that its learned noise variance is most beneficial when label corruption is systematic rather than uniform. This observation with VBLL is not alone; predictive performance on clean data is not well-correlated with noise robustness. In fact, some of the models with the lowest NDS across strategies such as NGBoost and SVM did not perform particularly well on clean data. These models may be more prone to picking up on noise rather than underlying signals in the data. 

% TODO: review, needs manual review
The ANOVA showed that model architecture and the model--representation interaction term dominate the variance in NDS. SVM and full BNNs stood out as noise-robust regardless of representation, and both rely on strong inductive biases, margin maximization in the case of SVM and weight priors in the case of BNNs. RF and NN-$\beta$ showed the opposite pattern, since they can be robust to noise but only when paired with particular representations, and these representations differ between the two models.

% TODO: review, needs manual review
Although representation is not a dominant determinant of NDS variance, it has an effect through the model--representation interaction term. On SMILES, model choice explains over 91\% of the variance in robustness, compared with 72\% on PDV. Across all six noise strategies, PDV configurations produced the shallowest mean NDS of all representations, with the best values coming from models with strong inductive biases such as SVM and full BNNs. This suggests that PDVs are preferable when dealing with noisy data, and it appears that the compact feature space of PDV complements margin-based and Bayesian regularization. Embeddings performed poorly, particularly when paired with NNs. These patterns are consistent with the importance of the interaction term and indicate that the choice of representation modulates noise robustness within a given architecture.

Taking a closer look at model choice for noise robustness, we ranked all model--representation configurations on QM9 by their Gaussian NDS. Tree-based ensembles dominate the top 10: NGBoost appears five times, XGBoost and RF twice each, and LightGBM once. No neural-network configuration appears the top 10. Rankings are consistent across strategies and the top QM9 configurations maintain their advantage on absorption, distribution, metabolism and excretion (ADME) datasets, though SVM configurations rank higher on these ADME datasets than NGBoost.

\section{Validation on Assay Datasets}
% [notes moved to thesis_notes_moved.md]

\begin{figure*}[htbp]
\centering
\includegraphics[width=\textwidth]{fig_validation_overview.png}
\caption[Noise degradation slope on LogD, Caco-2 and hERG $K_i$.]{Noise degradation slope (NDS) of five tree-based models under six noise strategies on the three assay datasets, LogD (top panel, a), Caco-2 efflux (middle panel, b) and hERG $K_i$ (bottom panel, c). The models are the random forest (RF), the quantile regression forest (QRF), XGBoost, LightGBM and NGBoost. More negative values (darker orange) indicate faster loss of performance, and the three panels share one color scale.}
\label{fig:validation_overview}
\end{figure*}

In addition to working with ``clean'' quantum molecular properties, we tested our findings on experimentally-obtained data, replicating our process of adding artificial noise. Unlike QM9, where $\sigma = 0$ corresponds to a noise-free baseline, the experimental endpoints carry their own measurement noise in both training and test labels; the injected $\sigma$ therefore stacks on top of an unknown noise floor. Per-dataset NDS heatmaps across all model architectures and noise strategies are shown in Figure~\ref{fig:validation_overview}. The ANOVA on these additional datasets confirms that model architecture dominates robustness variance on all three datasets, as seen in Additional file~\ref{app:noise:sec:assay_variance_decomposition}, and that the trends with respect to noise robustness generalize. NGBoost ranks first on both QM9 and external datasets, SVM also transfer quite well. Some models' predictive capabilities are limited on these external datasets, likely due to a combination of smaller training sets, narrower chemical coverage, and the unknown experimental noise in the labels. XGBoost and BNN variants suffer the most here, as seen in Figure~\ref{fig:validation_combined}a.

\begin{figure*}[htbp]
\centering
\includegraphics[width=\textwidth]{fig_validation_combined.png}
\caption[Noise robustness on QM9 and three assay datasets.]{Noise degradation slope (NDS) on QM9 (HOMO--LUMO gap, $N = 10{,}000$) and three assay datasets, LogD ($N = 5{,}039$), Caco-2 efflux ($N = 2{,}161$) and hERG $K_i$ ($N = 1{,}482$). The top panel (a) shows the NDS of seven models on each dataset (indicated by colors), where more negative values indicate faster loss of performance. The bottom panel (b) plots, for each model, the mean NDS over the three assay datasets against the mean NDS on QM9. RF = random forest, QRF = quantile regression forest, SVM = support vector machine, NN = neural network.}
\label{fig:validation_combined}
\end{figure*}

QRF was consistently less robust than RF on every external data set (Additional file~\ref{app:noise:sec:rf_vs_qrf_assay}). As seen with QM9, choice of molecular representation has a minimal effect on noise robustness across all three external datasets, reinforcing the finding that model architecture, not representation, is the main driver of noise robustness. 

\section{Uncertainty under Label Noise}\label{sec:uq}
% [notes moved to thesis_notes_moved.md]
% TODO: Hirschfeld2020, Scalia2020 and Norinder2014 are in refs.bib but not cited anywhere in the thesis. Cite them in this section or drop them.

The sections above showcased the importance of model selection when working with noisy labels. We take this one step further and ask, for the models which produce uncertainty estimates, how those estimates are affected by noise and if they can be used to track noise. 

\subsection{Per-Sample Uncertainty and Label Noise}
% [notes moved to thesis_notes_moved.md]

\citet{Kolmar2021} found that the mean predicted uncertainty derived from the GPs increases with the amount of label noise in the training data. We go beyond the population level and instead ask if a model's per-sample uncertainty tracks label noise. One key sign that a model is able to handle noise is its ability to track it, represented by the per-sample Spearman correlation between predicted uncertainty and noise magnitude (Table~\ref{tab:top_unc_noise}). While GPs and NGBoost gave the strongest correlations, BNNs were moderate. Learned embeddings (MHG-GNN and Mol2vec) struggled, so much so that their correlations dropped to near zero or below on every model tested. 

\begin{table}[htbp]
\centering
\caption[Strongest and weakest uncertainty--noise correlations on QM9.]{Model--representation combinations with the strongest and weakest correlations between predicted uncertainty and injected noise on the QM9 HOMO--LUMO gap under Gaussian noise. Unc-Noise $\rho$ is the Spearman correlation between predicted uncertainty and the magnitude of the injected noise $|\epsilon|$, and Unc-Error $\rho$ is the Spearman correlation between predicted uncertainty and the absolute prediction error. ECE is the expected calibration error, for which lower values indicate better calibration, and Cov.\ $1\sigma$ and Cov.\ $2\sigma$ are the empirical coverages of the $1\sigma$ and $2\sigma$ intervals, whose nominal values are 68\% and 95\%. The upper block lists the combinations with the highest Unc-Noise $\rho$ and the lower block those with the lowest. Bold marks the highest Unc-Noise $\rho$ for the GP and for NGBoost, and the best Unc-Error $\rho$ and ECE. GP = Gaussian process, BNN = Bayesian neural network, VBLL = Variational Bayesian Last Layer, SNS = Sort \& Slice, ECFP4 = extended-connectivity fingerprint of diameter four, PDV = physicochemical descriptor vector, MHG-GNN = Molecular Hypergraph Grammar graph neural network.}
\label{tab:top_unc_noise}
\small
\begin{adjustbox}{max width=\textwidth}
\begin{tabular}{llrrrrr}
\toprule
\textbf{Model} & \textbf{Rep.} & \textbf{Unc-Noise $\rho$} & \textbf{Unc-Error $\rho$} & \textbf{ECE} & \textbf{Cov. $1\sigma$} & \textbf{Cov. $2\sigma$} \\
\midrule
\textbf{GP}    & \textbf{SNS}    & \textbf{0.56} & 0.23 & 0.14 & 71.2\% & 94.6\% \\
GP             & ECFP4           & 0.53 & 0.22 & 0.14 & 71.4\% & 95.4\% \\
\textbf{NGBoost} & \textbf{SNS}  & \textbf{0.47} & 0.16 & 0.13 & 69.3\% & 95.4\% \\
\textbf{NGBoost} & \textbf{PDV}  & \textbf{0.47} & \textbf{0.29} & \textbf{0.11} & 69.0\% & 95.2\% \\
NGBoost        & PDV (binary)    & 0.44 & 0.25 & 0.13 & 69.8\% & 95.4\% \\
NGBoost        & ECFP4           & 0.44 & 0.20 & 0.14 & 69.8\% & 95.3\% \\
BNN-$\alpha$   & SNS             & 0.42 & 0.21 & 0.17 & 73.4\% & 95.1\% \\
BNN-$\beta$    & SNS             & 0.40 & 0.21 & 0.18 & 73.5\% & 95.1\% \\
BNN-$\beta$    & ECFP4           & 0.38 & 0.20 & 0.18 & 72.1\% & 95.0\% \\
BNN-$\beta$    & PDV             & 0.38 & 0.24 & 0.19 & 72.1\% & 94.0\% \\
\midrule
\multicolumn{7}{l}{\textit{Contrast: lowest correlations (graph/embedding representations)}} \\
\midrule
BNN-$\beta$ & Mol2vec & $-0.05$ & 0.02 & 1.33 & 50.6\% & 76.4\% \\
VBLL-$\alpha$ & MHG-GNN & $-0.05$ & 0.09 & 13.41 & 26.6\% & 38.6\% \\
VBLL-$\beta$ & MHG-GNN & $-0.03$ & 0.10 & 15.34 & 31.1\% & 47.1\% \\
VBLL-$\beta$ & Mol2vec & $-0.03$ & 0.15 & 23.86 & 44.8\% & 58.0\% \\
\bottomrule
\end{tabular}
\end{adjustbox}
\end{table}

\subsection{Treatment of Large Residuals}
% [notes moved to thesis_notes_moved.md]

Samples containing label noise tend to produce large residuals in training, and every model has its own way of handling them. NGBoost contains predicted scales that increase to absorb residual \citep{Duan2020}. The GP contains has a single global observation noise term. It does not absorb individual residuals, so its per-sample uncertainty derives primarily from the posterior variance rather than label noise  \citep{Rasmussen2005, Obrezanova2007}. BNNs treat these large residuals differently; they broaden their weight posteriors, increasing predicted uncertainty without explicitly modeling observation noise \citep{gal2016, kendall2017}. VBLL adds a learned noise variance to the BNN loss \citep{Harrison2024}, but this value is a global scalar, independent of the input. As seen in Table~\ref{tab:top_unc_noise} and Additional file~\ref{app:noise:sec:uncertainty_metrics_ecfp4}, both BNNs and VBLLs achieve moderate correlations which improve when paired with fingerprint or descriptor-based representations. Among models that learn separate distributional parameters during training, GP and NGBoost produced the strongest correlations between uncertainty and label noise, though QRF, which separates parameters using quantile distributions, produced both poor label predictions and uncertainty-noise correlations. At the population level, mean predicted uncertainty increases with artificial noise. Although we expect to see aleatoric uncertainty increase with injected noise, for VBLL both the aleatoric and epistemic components increased (Figure~\ref{fig:uncertainty_combined}b).

\begin{figure}[htbp]
\centering
\includegraphics[width=\textwidth]{fig_uncertainty_combined.png}
\caption[Mean predicted uncertainty under label noise on QM9.]{Mean predicted uncertainty against the injected noise level $\sigma$ on the QM9 HOMO--LUMO gap. The top panel (a) shows the total uncertainty of six probabilistic models, namely the quantile regression forest (QRF), NGBoost, and the full Bayesian (BNN) and Variational Bayesian Last Layer (VBLL) variants of NN-$\alpha$ and NN-$\beta$. The bottom panel (b) splits the uncertainty of VBLL-$\alpha$ and VBLL-$\beta$ into aleatoric (solid lines) and epistemic (dashed lines) components.}
\label{fig:uncertainty_combined}
\end{figure}

\subsection{Aleatoric and Epistemic Attribution}
% ---- NOTES: Uncertainty attribution is itself representation-dependent ----
% Takeaway: the variance-head BNNs put added noise into the aleatoric part, as expected. QRF puts much of it into the epistemic part, and the representation decides how much.
% Takeaway: the aleatoric part shows that noise rose, not how much. Under censoring, uncertainty shrinks instead of rising.
% Paragraph 1: the question, and which models can answer it.
% - Mean uncertainty rises with noise, but that does not say what the model attributes the rise to.
% - Expected behavior: the aleatoric part rises with injected label noise; the epistemic part stays roughly flat.
% - Three model kinds give a per-molecule split: the full BNNs with a variance head, the heteroscedastic GP, and QRF \citep{kendall2017, Rasmussen2005, Meinshausen2006}.
% - The other GPs are excluded: their observation noise is one number for the whole fit, so it cannot follow individual molecules \citep{Rasmussen2005}.
% - Measure: the rise in mean epistemic uncertainty divided by the rise in mean aleatoric uncertainty, from no added noise to noise level 1.5, over held-out molecules.
% - A ratio near 0 means the added noise went to the aleatoric part. A ratio of 1 or more means the epistemic part grew as much or more.
% Paragraph 2: the variance-head BNNs split as expected; QRF does not.
% - Full BNNs with a variance head: the ratio is at most 0.43 in every model, dataset, representation and noise condition.
% - QRF: both parts rise with noise, and on some representations the epistemic part rises faster.
% - This holds although QRF is among the more noise-robust models: on QM9 under Gaussian noise it ranks 4th of the 19 models and variants on AUC_norm on four of the six representations.
% - Do not call QRF among the most accurate. On the same data its clean R^2 ranks between 6th and 13th of the 19, depending on the representation.
% Paragraph 3: for QRF the representation decides whether the split works.
% - Contrast: representation explained little of the variance in robustness, but here it decides the split.
% - QRF ratio on held-out molecules, range over datasets and every noise condition except censoring: 0.63 to 1.09 on the three fingerprints (ECFP4, Avalon, Sort & Slice).
% - Same range on PDV: 1.03 to 1.42. On ChemBERTa and MHG-GNN: 1.28 to 2.05.
% - Use the held-out range for PDV (1.03 to 1.42), not the paper's 1.01 to 1.45, which pools held-out and out-of-fold training molecules.
% - Under Gaussian noise alone, held-out molecules: ECFP4 0.65-0.75, Sort & Slice 0.69-0.76, Avalon 0.81-0.89, PDV 1.17-1.32, MHG-GNN 1.54-1.72, ChemBERTa 1.60-1.88.
% - So on the learned embeddings QRF puts more of the added noise into its epistemic part than into its aleatoric part.
% Paragraph 4: grouped-shifted noise moves the heteroscedastic GP, most on fingerprints.
% - Heteroscedastic GP, Gaussian noise, held-out molecules: ratio between -0.05 and 0.36 on every representation.
% - Heteroscedastic GP, grouped-shifted, range over datasets: ECFP4 0.73-1.00, Sort & Slice 0.44-0.61, Avalon 0.32-0.46.
% - Same model and condition: PDV 0.13-0.22, ChemBERTa 0.14-0.23, MHG-GNN 0.21-0.25.
% - Laplace, Student-t and outlier noise give splits like those under Gaussian noise.
% - Grouped-shifted raises the epistemic share of the heteroscedastic GP and of QRF for every dataset and representation.
% - Possible reading, stated as a reading: a shared offset per scaffold family looks to the model like structure in the data, not label noise.
% Paragraph 5: the aleatoric part shows that noise rose, not how much.
% - This comparison is against the injected noise, so it uses out-of-fold training molecules. Held-out labels carry no injected noise.
% - Slope of mean aleatoric uncertainty against the standard deviation of the added noise, both in the label's units, Gaussian noise, median over folds.
% - On ECFP4 and PDV: 0.36 (QRF, PDV, Caco-2) to 1.09 (a full BNN with a variance head, PDV, hERG K_i).
% - Change: the paper's text says 0.36 to 0.99. The upper end from the fitted slopes is 1.09; re-check the paper's source.
% - Including ChemBERTa, the range widens to 0.28 (QRF, Caco-2) to 1.62 (a full BNN with a variance head, hERG K_i).
% - So it can compare noise between datasets or noise levels. It does not estimate a dataset's own measurement error.
% Paragraph 6: censoring shrinks the uncertainty instead of raising it.
% - Clipping labels at a limit narrows their spread.
% - NGBoost and QRF shrink their total uncertainty with that spread on every dataset; the GPs do the same on hERG K_i and Caco-2.
% - None of these models treats a clipped label as a noisy one, so they grow more confident while losing information.
% Table: ratio of epistemic rise to aleatoric rise. One row per model (the two full BNNs with a variance head, heteroscedastic GP, QRF) and noise condition (Gaussian, grouped-shifted).
% - One column per representation (ECFP4, Avalon, Sort & Slice, PDV, ChemBERTa, MHG-GNN). Each entry is the range over datasets, held-out molecules.
% Figure: aleatoric and epistemic parts against noise level, one panel per model with a per-molecule split, QM9 HOMO--LUMO gap, ECFP4, Gaussian noise.
% Figure: aleatoric part against the standard deviation of the added noise, per model, ECFP4 and PDV, out-of-fold training molecules, with a dashed line where reported noise equals added noise.
% ---- END NOTES ----

% CODE INPUT
Mean predicted uncertainty rises with label noise, but this alone does not show to which component a model attributes the rise. We expected the aleatoric component to rise with injected label noise and the epistemic component to stay roughly flat. Three kinds of model split their uncertainty per molecule, namely the full BNNs with a variance head, the heteroscedastic GP and QRF \citep{kendall2017, Rasmussen2005, Meinshausen2006}. We excluded the other GPs because their observation noise is a single number for the whole fit and therefore cannot follow individual molecules \citep{Rasmussen2005}. For each model, we divided the rise in mean epistemic uncertainty on held-out molecules, from no added noise to a noise level of 1.5, by the rise in mean aleatoric uncertainty over the same range. A ratio near 0 means that the added noise went to the aleatoric component, while a ratio of 1 or more means that the epistemic component grew as much or more (\textbf{[FIGURE NEEDED: table of the ratio of epistemic rise to aleatoric rise, one row per model (the two full BNNs with a variance head, heteroscedastic GP, QRF) and noise condition (Gaussian, grouped-shifted), one column per representation, each entry the range over datasets on held-out molecules]}).

% CODE INPUT
The full BNNs with a variance head split the added noise as expected, with a ratio of at most 0.43 for every model, dataset, representation and noise condition. QRF behaved differently, since both of its components rise with noise and on some representations the epistemic component rises faster (\textbf{[FIGURE NEEDED: aleatoric and epistemic parts against noise level, one panel per model with a per-molecule split, QM9 HOMO--LUMO gap, ECFP4, Gaussian noise]}). This holds although QRF is among the more noise-robust models, ranking 4th of the 19 models and variants on \aucnorm on four of the six representations on QM9 under Gaussian noise. Its clean $R^2$ on the same data ranks between 6th and 13th of the 19, depending on the representation.

% CODE INPUT
Representation explained little of the variance in noise robustness, but for QRF it determines how the added noise is split. Over all datasets and every noise condition except censoring, the QRF ratio on held-out molecules ranges from 0.63 to 1.09 on the three fingerprints (ECFP4, Avalon and Sort \& Slice), from 1.03 to 1.42 on PDV, and from 1.28 to 2.05 on ChemBERTa and MHG-GNN. Under Gaussian noise alone, the held-out ranges are 0.65--0.75 on ECFP4, 0.69--0.76 on Sort \& Slice and 0.81--0.89 on Avalon, compared with 1.17--1.32 on PDV, 1.54--1.72 on MHG-GNN and 1.60--1.88 on ChemBERTa. On the learned embeddings, QRF therefore attributes more of the added noise to its epistemic component than to its aleatoric component.

% CODE INPUT
The type of noise also shifts the split, most clearly for the heteroscedastic GP. Under Gaussian noise, its ratio on held-out molecules lies between $-0.05$ and 0.36 on every representation. Under grouped-shifted noise, the range over datasets is 0.73--1.00 on ECFP4, 0.44--0.61 on Sort \& Slice and 0.32--0.46 on Avalon, compared with 0.13--0.22, 0.14--0.23 and 0.21--0.25 on PDV, ChemBERTa and MHG-GNN, respectively. Laplace, Student-t and outlier noise give splits similar to those under Gaussian noise. Grouped-shifted noise raises the epistemic share of both the heteroscedastic GP and QRF for every dataset and representation. One possible explanation is that a shared offset within each scaffold family appears to the model as structure in the data rather than as label noise.

% CODE INPUT
We next asked how closely the aleatoric component follows the amount of added noise. Because held-out labels carry no injected noise, this comparison uses out-of-fold training molecules. We fitted the slope of mean aleatoric uncertainty against the standard deviation of the added noise, both in the units of the label, under Gaussian noise, and we report the median over folds (\textbf{[FIGURE NEEDED: aleatoric part against the standard deviation of the added noise, per model, ECFP4 and PDV, out-of-fold training molecules, dashed line where reported noise equals added noise]}). On ECFP4 and PDV, the slope ranges from 0.36, for QRF on PDV for Caco-2, to 1.09, for a full BNN with a variance head on PDV for hERG $K_i$. Including ChemBERTa widens the range to between 0.28, for QRF on Caco-2, and 1.62, for a full BNN with a variance head on hERG $K_i$. The aleatoric component can therefore be used to compare noise between datasets or between noise levels, but it does not estimate the measurement error of a dataset.

% CODE INPUT
Censoring behaves differently from the other noise conditions, because clipping labels at a limit narrows their spread. NGBoost and QRF shrink their total uncertainty with that spread on every dataset, and the GPs do the same on hERG $K_i$ and Caco-2. None of these models treats a clipped label as a noisy one, so they become more confident while the labels lose information.

% TODO: MISSING RESULT: the GP uncertainty decomposition by representation.

\subsection{Representation and Uncertainty Tracking}
% [notes moved to thesis_notes_moved.md]

The choice of representation has a much stronger impact on per-sample uncertainty tracking than on noise robustness. Fingerprints and physicochemical descriptors provided a clear signal with which models such as GP and NGBoost could separate noise-induced residuals from structural variation. However, learned embeddings such as MHG-GNN and Mol2vec did not produce viable uncertainty-noise correlations (Table~\ref{tab:top_unc_noise}). Although these learned embeddings can be effective representations for many tasks, this research suggests that they provide challenges for per-sample noise tracking. However, there is room for further investigation, as only two embeddings were tested.

\section{Conclusions}
% [notes moved to thesis_notes_moved.md]

These results lead to two main conclusions. The first is that noise robustness, defined as the rate at which predictive performance degrades with increasing label noise, is primarily determined by the model's training mechanism. The margin maximization of SVM, the weight priors of BNNs, and the ensembling in tree methods all appear to keep those models from fitting to noise. The choice of representation has a stronger influence on predictive performance than noise robustness. The interaction term between representation and model has a stronger impact; some models like SVM and full BNNs remain robust across all representations, while RF and NNs are robust only with certain ones. While the pattern of noise has a strong impact on predictive performance, it does not influence noise robustness rankings between models or representations much.

We investigated how a wide variety of molecular representations and model architectures perform under increasing label noise in QSAR modeling. We used ANOVA variance decomposition across three noise conditions on the QM9 HOMO--LUMO gap task to show that with regards to predictive performance, model architecture and molecular representation account for comparable shares of the variance. However, for robustness to label noise, measured as \aucnorm, the normalized area under the $R^2$ retention curve, model architecture is the dominant factor, while representation explains less than 10\% of the variance.

Threshold and value-proportional noise produce the widest spread of NDS across models; outlier noise barely separates them. Model rankings are highly concordant across strategies (Kendall's $W = 0.92$), suggesting that a model's robustness to noise is independent of type of noise. Predictive performance on clean data is not correlated with noise robustness; NGBoost and SVM, the most noise-robust models, did not perform particularly well on clean data. Embeddings (MHG-GNN, Mol2vec) struggled with regards to noise robustness, especially with neural network architectures. On the flip side, QRF was significantly less noise-robust than RF across datasets, suggesting that quantile regression may be overfitting to noisy labels rather than absorbing them.

% TODO: review, needs manual review
The second conclusion concerns whether the uncertainty estimates of a model track per-sample label noise. Some models struggle to do so, particularly BNNs, whose uncertainty comes from a posterior over weights, and QRF, which produces a full quantile distribution but does not track the noise successfully. Models that fit a separate per-sample scale or observation-noise parameter during training track label noise more effectively. NGBoost identifies noisy labels through a per-sample predicted scale that absorbs large residuals. The GP does not, because it has a single global observation-noise term. Its per-sample uncertainty therefore comes from the posterior variance, which depends on the position of a molecule relative to the training molecules and not on its label. Among representations, uncertainty tracking collapses with the learned embeddings MHG-GNN and Mol2vec, while fingerprints and descriptors produce a usable signal. Neither strong predictive performance nor noise robustness implies noise tracking, which again depends on the architecture of the model and on how it handles the loss. 

Although the choice of representation carries less weight than that of models, we observed an interesting phenomenon. PDVs produce the strongest predictive performance on clean data, but are the least noise-resistant. Embeddings, particularly Mol2vec, are the most robust to noise yet perform the worst noise tracking. At the end of the day, the choice of both model and representation comes down to your problem, including the amount of noise in your data, your objectives, and your compute limits. 

QM9 served as a relatively clean baseline; it is computationally-derived with negligible measurement noise. The validation datasets, LogD, Caco-2 \citep{openadmet}, and hERG $K_i$ \citep{Zdrazil2023}, tested are smaller and cover a narrow range of the chemical and biological space, and have an unknown inherent amount of noise from experimentally-determined values, on top of the added artificial noise. These validation datasets supported the QM9 findings, with XGBoost the notable exception, degrading on the external datasets. Future work should extend this type of noise robustness benchmarking to classification tasks and larger experimental datasets with better-known uncertainty estimates.

When working with noisy experimental data in QSAR settings, it is important to keep in mind that noise robustness and noise-aware uncertainty estimation are primarily driven by the model's mechanism of training, and to a certain extent the choice of molecular representation, usually with respect to the model. While Bayesian transformations on neural networks improve noise robustness, they do not produce uncertainty estimates that correlate particularly strongly with noise. Models like NGBoost and GPs, which include a learned noise or scale parameter in the training objective, are the most robust to noise across both various types of noise and datasets and often produce uncertainty estimates which track per-sample label noise. Pairing these models with fingerprint-based representations provides additional benefits in restricting, and potentially detecting and mitigating label noise. 

